Priorytet określa ważność funkcji dla użytkownika i wynikające z niej ryzyko produktowe,

które wyznacza intensywność testowania (testowanie oparte na ryzyku).

{\small
\begin{longtable}{p{0.188\linewidth}p{0.188\linewidth}p{0.188\linewidth}p{0.188\linewidth}p{0.188\linewidth}}
\toprule
\textbf{ID} & \textbf{Funkcja} & \textbf{Opis} & \textbf{MoSCoW} & \textbf{Uzasadnienie} \\
\midrule
\endhead
F01 & Rejestracja konta & Założenie konta przez formularz \texttt{/user/sign\_up} z walidacją nazwy, e-maila i hasła (min. 8 znaków). & Must & Wejście każdego nowego użytkownika do systemu. \\
F02 & Logowanie i wylogowanie & Uwierzytelnienie nazwą/e-mailem i hasłem, sesja w ciasteczku, wylogowanie z menu użytkownika. & Must & Warunek wstępny niemal wszystkich pozostałych funkcji. \\
F03 & Zarządzanie repozytoriami & Tworzenie (publiczne/prywatne, inicjalizacja README), edycja ustawień i usuwanie repozytoriów. & Must & Podstawowy obiekt domeny aplikacji. \\
F04 & Przeglądanie kodu i historii & Lista plików, podgląd README, historia commitów, gałęzie. & Must & Główny sposób korzystania z repozytorium. \\
F05 & Zgłoszenia (issues) & Tworzenie, komentowanie, zamykanie i ponowne otwieranie zgłoszeń. & Must & Kluczowy proces śledzenia błędów i zadań. \\
F06 & Pull requesty & Porównanie gałęzi, utworzenie PR, przegląd zmian, merge. & Must & Rdzeń współpracy nad kodem. \\
F07 & Kontrola dostępu & Widoczność repo prywatnych, kolaboratorzy z uprawnieniami, ochrona zasobów przed anonimami. & Must & Poufność danych – naruszenie ma wysoki wpływ. \\
F08 & Edytor plików w przeglądarce & Edycja pliku w UI i commit bezpośrednio do gałęzi lub nowej gałęzi. & Should & Wygodne, ale istnieje alternatywa (klient git). \\
F09 & Etykiety i kamienie milowe & Kategoryzacja zgłoszeń etykietami, grupowanie w milestone z licznikiem postępu. & Should & Organizacja pracy zespołu. \\
F10 & Wyszukiwanie & Wyszukiwanie repozytoriów, użytkowników i zgłoszeń (Explore). & Should & Istotne przy dużej liczbie repozytoriów. \\
F11 & Organizacje i zespoły & Grupowanie repo w organizacjach, zespoły z uprawnieniami. & Could & Potrzebne większym zespołom; poza krytyczną ścieżką. \\
F12 & Profil i tokeny API & Ustawienia konta, tokeny dostępu do REST API z zakresami (scopes). & Could & Wykorzystywane przez integracje i automatyzację. \\
F13 & Wiki repozytorium & Strony dokumentacji przy repozytorium. & Could & Funkcja dodatkowa. \\
F14 & Gitea Actions (CI) & Wbudowany system CI kompatybilny z GitHub Actions. & Won't & Poza zakresem – wymaga runnerów, osobny przedmiot testów. \\
\bottomrule
\end{longtable}
}
Podsumowanie: 7 $\times$ Must, 3 $\times$ Should, 3 $\times$ Could, 1 $\times$ Won't.

Funkcje Must są pokryte przypadkami funkcjonalnymi na poziomie UI i/lub API oraz testami

niefunkcjonalnymi; F11 i F13 świadomie nie są pokryte (ryzyko rezydualne, zob. raport §8).

